% ============================================================
% AUTOMATISCH GEGENEREERD BESTAND
% Niet handmatig aanpassen.
% ============================================================

\ifdefined\HCode
  % Geen afzonderlijke module-openingspagina in HTML.
\else
  \FVDmoduleopening
    {Van krachten tot kwanta}
    {In deze cursus bouwen we verder op de fysica uit het vijfde jaar en onderzoeken we hoe krachten, energie, trillingen, golven en licht met elkaar samenhangen. We starten vanuit de klassieke mechanica. We bestuderen krachten, evenwicht, arbeid en energie en onderzoeken hoe fysische systemen beschreven en voorspeld kunnen worden met behulp van modellen en wiskundige verbanden. Vervolgens richten we ons op trillingen en golven. Vanuit harmonische trillingen bouwen we stap voor stap het golfmodel op. We onderzoeken hoe golven zich voortplanten en bestuderen verschijnselen zoals reflectie, breking, interferentie, diffractie, resonantie en het Dopplereffect. Geluid vormt daarbij een belangrijke toepassing. Daarna passen we het golfmodel toe op licht. We onderzoeken licht zowel met het stralenmodel als met het golfmodel en bestuderen onder andere reflectie, breking, lenzen, beeldvorming en elektromagnetische golven. Ten slotte verkennen we de grenzen van de klassieke fysica. Aan de hand van enkele ideeën uit de moderne fysica onderzoeken we hoe concepten uit de kwantumfysica en de relativiteitstheorie ons beeld van materie, straling, ruimte en tijd hebben veranderd. Doorheen de cursus gebruiken we wiskunde als taal om fysische verschijnselen te beschrijven en te modelleren. Vectoren, functies, goniometrie, afgeleiden en integralen krijgen zo een concrete fysische betekenis. Experimenten, gegevensanalyse en probleemoplossend denken helpen ons om modellen kritisch te toetsen aan de werkelijkheid. Fysica wordt zo geen verzameling afzonderlijke formules, maar een samenhangend geheel van ideeën waarmee we de natuur kunnen beschrijven, verklaren en voorspellen.}
    {%
    \FVDmodulethemetile{1}{Beweging en kracht — Van beweging beschrijven tot energie begrijpen}
    \FVDmodulethemetile{2}{Alles trilt — Van trilling tot golf}
    \FVDmodulethemetile{3}{Licht in beeld — Van lichtstraal tot elektromagnetische golf}
    \FVDmodulethemetile{4}{Voorbij het klassieke — Van klassieke fysica tot kwanta en relativiteit}
    }
\fi
